\begin{abstract}
对未知空间进行建图，并不必然揭示另一台机器人能够完成操作的位置。
本文研究这一矛盾在空地协同取回任务中的表现，提出由地面操作需求引导空中观测的主动感知框架。
该框架将经过验证的逆可达性候选，通过地面机器人精确位姿处的底盘占地轮廓投影到环境空间，
识别与后续操作相关的环境证据。在有限时间激光扫描模型中，这一支撑对观测缺失度进行加权，
用于选择无人机的下一观测视点。实测地面支撑和考虑碰撞的候选筛选将感知目标连接到物理执行；
预测观测不替代真实测量。
在由 P450、BUNKER、AUBO i5 和 AG95 组成的物理仿真平台上，本文使用 96 个独立场景、
192 次主配对任务验证完整闭环。在相同的三个观测窗口预算下，本文方法（Ours）在 80/96 个场景中
实现物理取回，而 Generic NBV 为 56/96。两种策略共用候选集、传感器模型、成本与执行链，
仅观测增益的空间加权不同。配对成功率差为 25.00 个百分点，预先规定的保守 95\% 置信区间为
8.56--39.15 个百分点，双侧精确 McNemar 检验 $p=8.43\times10^{-6}$。
这些结果支持在静态、已知物体的取回任务中利用后续操作需求分配空中感知资源，
但尚未确立普遍的时间节省或实机性能优势。
\end{abstract}
